\section{Middle Earth Once More}

\begin{quotex}
By way of such aspects, we see---it is clear---in the Middle Ages an awakening of the true forces already acting in Nordic-Aryan Romanity, of its true solarity, propitiated from such a resurgence or rebirth from a new contribution of Aryan blood...

\flright{Gornahoor translating Evola\footnote{\url{https://gornahoor.net/?p=3861}}, on Cesaro}
\end{quotex}

Indeed, a \emph{tremendously sophisticated} defense of the very opposite view (which Evola here castigates) can be found by a modern feminist\footnote{\url{http://www.bu.edu/arion/files/2010/03/Paglia-Great-Mother1.pdf}}, which draws on both Neumann and J. Bachofen in truly classical fashion, in order to defend the earth-cults and the ``rise of the feminine''. Yet it should be obvious by now to all those with perception that the ``rise of feminism'' entails the non-existence (in experience) of the solar origins, whereas the opposite is not true. Solarity can comprehend and order what is below it, but what is below it cannot comprehend and order solarity -- this is the meaning (of course also) of the Biblical myth and mandate for the man to be the spiritual ``head'' of the woman, although it is typically taken very literally. Thus, ``what comes after'' is not a novelty or evolution, but rather a devolution. This is exactly the opposite of that which is generally ``held'' in the academic and spiritual worlds of today, where progressives try to convince others that they have ``something better to offer'' in tune with the Zeitgeist than those horrible old days of long ago when ``we'' (who is this ``We'', here?) didn't know any better. This is the world of the Iron Polygon\footnote{\url{http://unqualified-reservations.blogspot.com/2007/05/iron-polygon-power-in-united-states.html}} and the ``Cathedral\footnote{\url{http://unqualified-reservations.blogspot.com/2009/01/gentle-introduction-to-unqualified.html}}'', where everyone has a ``right'' to ``rights'', which usually means squabbling over personally disgusting political material. Some might argue for the ``Matrix'' as a better word here, and I am sure the terminology and word play will continue to be plied, in order to come to a more perfect understanding of how the Modern World is evolving.\\
{}\strut \\

But ``devolution'' is a tremendous theme that recurs over and over again with Evola. An acquaintance put it this way:

\begin{quotex}
(1) Evola dealt with both the seen and the unseen, the commonly recognized and perceivable, and the obscure.

(2) Evola dealt with the controversial --- the things which salaried men dependent upon their peers' approval refuse to speak on.

(3) Evola represents a regal framework which recognizes the proper stance of a solider, the transcendental outlook. This is also imperial insofar as the imperium stretches beyond any exception that a man or people might make for itself.

(4) Evola, even if his wartime endeavors may have been misguided, was always forward looking. His thought is relevant to the present day.

(5) Because of the above, Evola thought, in uncommon fashion, whereas most others simply regurgitate the opinions of their peers.

(6) Evola was one of a few universal minds, both pan-European and looking seriously at the cultures of the East, of the type that are no longer found today among a worldwide culture which no longer takes the traditional arts seriously and tears down those who do.

(7) Evola plotted a path for individuals to overcome, to transcend their individual limitations, without dependence on dogma.

(8) Evola did not mince words, neither adding fluff in order to sell himself to the masses nor adding mean witticisms for the sake of producing petty laughter.

(9) Evola was not motivated by fear, positing or selling post-mortem metaphysical realms for those whose hearts are too faint to face the present one. He could not and did not divorce the political and metaphysical --- the danger of all priests who implicitly attempt to play at politics.

(10) Evola was rooted and established in the transcendental, the spiritual man spoken of and lauded by the ancients, the sage who is known to those who know they do not understand.
\end{quotex}


(Credit to Joel Dietz, of The Conservatory, Dunedain.net)

Steiner's idea of ``Evolution'' can (ultimately) be said to de-rail his entire system, no matter what his other virtues and how much other truth he contains, or personally embodied. This could account for the death of his movement after his passing, for there was no doubt that Steiner knew a great deal of mysterious and occult things. Evola, situated as he was, was in a better place to understand how little compromise was possible with the new conditions; he knew that the Revolution intended to dethrone the true Self in the name of the sensate individuality. Steiner, on this point, illustrates how devastating even a minor flaw in Reason can be, a point Alain Benoist is also illustrating before our eyes in his current career (as Gornahoor has also highlighted).

Like the Roman Catholic Church showed man, a great deal of experiential freedom can be tolerated, but not certain fundamental theoretical errors which go deep enough, leading (in the end and if pursued vigorously) to the ``triumph of man's reason\footnote{\url{http://en.wikipedia.org/wiki/Revolt_in_the_Vend\%C3\%A9e}}''.

There is a ``Christian'' reason for this very fact, as well, which is illustrated over and over again, countless times, in the annals of man (Evola learned quite well from his Catholic teachers). Errors do not immediately manifest their fruits (which is why they are ``fruits''), at least not to those of us in between heaven and earth. The only thing that can help here is to be ``wise as serpents'', or, (as Jesus illustrated) to use one's divinely inspired and lead Reason to steer through the narrow gate, for Reason (the Noetic faculty) is the only portion of man\footnote{\url{http://www.youtube.com/watch?v=qvfbocYAs7w}} which a ``beginner'' can use with impunity from his dark side.

The good thing is that there are ``fruits'' which can be held up to be judged. Were the Middle Ages really that bad\footnote{\url{http://listverse.com/2009/01/07/top-10-myths-about-the-middle-ages/}}? As credit substitutes for open slavery in late stage capitalism, can one continue to maintain that everything ancient regime is distasteful and misguided? It is inconceivable to the modern mindset that anything important has been lost, since it is governed by instrumental ``reason'', functionalism, and a pragmatism which is ideologically charged to rationalize whatever it cannot account for strictly speaking, on even its own terms. It is also precommitted to what Tomberg would call the purely ``horizontal'', or ``circle of the serpent'': we do not need God (or the higher Self) because we now have ``Infinity'' (or perhaps, ``Singularity''). Pages of narrative are necessary to qualify how this ``gutting'' of transcendence can preserve any meaning at all.

The Middle Ages (if one will actually open the annals\footnote{\url{http://www.flw.ugent.be/btng-rbhc/pdf/BTNG-RBHC,\%2034,\%202004,\%204,\%20pp\%20567-579.pdf}} and writings of history, or even go to a cathedral to study them\footnote{\url{http://www.gutenberg.org/ebooks/24428}}) are a permanent refutation of the idea and ideal of Progress. How can ``Progress'' mean anything when confronted with the undeniable fact that the Middle Ages, \emph{even on their own progressive theories}, bridge the Ancient collapse and the Modern Day? What will they refute a Gothic rose window with? They were the time of the ``truly human'', living in a precarious but cheerful balance (even Nietzsche saw their cheerfulness tinging as late a century as the 17th), between the age of the gods and the age of the animals. They were perfectly willing\footnote{\url{http://listverse.com/2007/09/22/top-10-inventions-of-the-middle-ages/}} to utilize technology, when a cultural space could be found for such, but were unwilling to commit themselves wholesale to a view of man as a ``plug-in'' to his own devisings.

As the Modern World deepens its devolution, the Middle Ages will continue to grow lighter and brighter. It will become more crystal clear that there is only way out -- not a path ``deeper into Chaos'' but a path back to where we have been told ``There Be Dragons''.

The world already awaits us.


\flrightit{Posted on 2012-03-16 by Logres}

